% === TABLE LEGENDS ===
% Table caption text, listed separately from the table float (MRM style, mirrors
% the Figure Legends). Number auto-synced via \ref to the table's \label.
%TC:ignore
\section*{Table Legends}

\noindent\textbf{Table~\ref{tab:auc}.}
Tumor-detection performance: the spectral classifier matches ADC, and the
two outer fractions alone suffice.
With one exception, entries are leave-one-out cross-validated AUCs (95\%
confidence interval) for tumor-versus-normal classification, evaluated through
an identical $L_2$-regularized logistic-regression pipeline ($C = 1.0$,
standardized features) so that single- and multi-feature models are directly
comparable (Figure~\ref{fig:roc}); the fair head-to-head is therefore ADC
(logistic regression) versus the spectral models.
The exception, ADC (raw rank), is the conventional in-sample ranking by ADC
value (no model fit), reported as a reference; it is the hardest baseline to
beat, yet the spectral models match it within overlapping confidence intervals.
``8 fractions'' uses all spectral bins; ``2 outer fractions'' uses only
$D = 0.25$ and $3.0\;\umsmm$.
Confidence intervals are percentile bootstrap over patients
(2000 resamples) \citep{efron1993bootstrap}.
Individual fractions are NUTS single-feature raw-rank AUCs---the per-fraction ROC
curves in Figure~\ref{fig:roc}.
Grade-classification AUCs are omitted: with only $n = 29$ graded tumors they are
too imprecise to be informative (95\% intervals spanning $\sim$0.5--0.95); the
grading analysis is instead presented as grade-dependent spectral shifts
(Figure~\ref{fig:spectrum_ggg}) and rank correlations (Results).
%TC:endignore
